% =========================================================================
% RESUMO ESTENDIDO EM LÍNGUA PORTUGUESA
% OBRIGATÓRIO PARA DISSERTAÇÕES E TESES REDIGIDAS EM LÍNGUA INGLESA
% (Conforme Resolução do CEPE/UnB e diretrizes institucionais do PPGT)
% =========================================================================
% INSTRUÇÕES:
% 1. Se a sua dissertação/tese for escrita em INGLÊS (\selectlanguage{english}),
%    a Universidade de Brasília exige a inclusão obrigatória de um Resumo
%    Estendido em português com extensão típica de 2 a 5 páginas.
% 2. O ambiente {otherlanguage}{brazilian} abaixo garante que a hifenização
%    e separação silábica deste texto em português permaneçam perfeitas,
%    mesmo com o documento principal configurado em língua inglesa.
% =========================================================================

\begin{otherlanguage}{brazilian}
\chapter*{RESUMO ESTENDIDO}
\label{resumoestendido}

Este é o modelo de \textbf{Resumo Estendido em Língua Portuguesa}. Segundo as normas regimentais da Universidade de Brasília (Resolução CEPE/UnB) e do Programa de Pós-Graduação em Transportes (PPGT), todo trabalho de conclusão de mestrado ou doutorado redigido em língua estrangeira (como o inglês) deve, obrigatoriamente, incluir um resumo aprofundado em língua portuguesa com extensão típica entre 2 e 5 páginas.

O texto do resumo estendido deve apresentar uma síntese clara e estruturada de toda a pesquisa, contemplando os tópicos essenciais detalhados a seguir.

\section*{Contextualização e Formulação do Problema}

Apresente a fundamentação geral do problema de transportes, sua relevância social, econômica, operacional ou ambiental, destacando os desafios enfrentados pela engenharia de transportes e a motivação do estudo.

\section*{Objetivos da Pesquisa}

Explicite com precisão o objetivo geral da pesquisa e os objetivos específicos estabelecidos para alcançá-lo.

\section*{Procedimentos Metodológicos}

Descreva a abordagem metodológica adotada (experimental, empírica, simulação computacional, modelagem matemática, econométrica ou pesquisa qualitativa), detalhando:
\begin{itemize}
    \item As bases de dados empregadas ou cenários de coleta de dados primários;
    \item Os modelos teóricos, algoritmos ou formulações conceituais aplicados;
    \item Os métodos de calibração, validação ou tratamento estatístico utilizados.
\end{itemize}

\section*{Principais Resultados e Discussão}

Sintetize os achados fundamentais obtidos a partir das análises, destacando métricas quantitativas, tendências observadas, comparações com a literatura especializada e o desempenho das intervenções ou modelos propostos.

\section*{Conclusões e Contribuições}

Apresente as conclusões centrais da pesquisa, ressaltando:
\begin{itemize}
    \item O atendimento aos objetivos propostos inicialmente;
    \item As contribuições científicas, metodológicas e práticas para o setor de transportes;
    \item As limitações identificadas no desenvolvimento do estudo;
    \item Sugestões concretas de pesquisas e desdobramentos futuros.
\end{itemize}

\end{otherlanguage}

